\section{Retraction of covering centres for preparation targets}\label{sec:centres70}
The centres in a covering problem need not belong to the family being
covered. For preparation targets this distinction can be removed exactly,
without a small-error or rank assumption. Throughout this section a centre
is one memoryless instrument, reused at every call; it is not a device with
an additional persistent environment.

\begin{theorem}[Preparation-centre retraction]\label{thm:retraction70}
Let $\mathcal C=(\mathcal C_y)_{y=1}^m$ be any instrument from $M_d$ to
$M_n$, and set
\begin{equation}\label{eq:retraction70}
 \tau_y=\mathcal C_y(I_d/d)=d^{-1}\operatorname{Tr}_{\rm in}J_y(\mathcal C),
 \qquad \mathfrak R\mathcal C=\mathcal R_\tau.
\end{equation}
Then $\mathfrak R$ maps legal instruments to preparation instruments,
fixes every preparation instrument, and preserves rational Choi data.
For every $N\ge1$ and every $\sigma\in\mathcal S_{m,n}$,
\begin{equation}\label{eq:centredominance70}
 d_N(\mathcal R_\sigma,\mathfrak R\mathcal C)
                  \le d_N(\mathcal R_\sigma,\mathcal C).
\end{equation}
Here $d_N$ is the unhalved trace distance optimized over the common
adaptive testers of Lemma~\ref{lem:productreplacer68}. Consequently, for
any subset of preparation targets and any radius $\delta>0$, the minimum
covering cardinality is unchanged when arbitrary legal memoryless centres
are restricted to preparation centres. The assertion also holds when
centres are required to have rational Choi entries.
\end{theorem}
\begin{proof}
Complete positivity gives $\tau_y\ge0$, and trace preservation gives
$\sum_y\tr\tau_y=1$. The formula is rational and is the identity on the
preparation family, so $\mathfrak R^2=\mathfrak R$. Feed the state $I_d/d$
into each of $N$ calls to $\mathcal C$, use a fresh input at the next call,
and retain all outputs and outcome labels. The resulting state is
$\boldsymbol\tau^{\otimes N}$, since the centre is memoryless. The same
tester applied to $\mathcal R_\sigma$ returns
$\boldsymbol\sigma^{\otimes N}$. Therefore
\[
 \|\boldsymbol\sigma^{\otimes N}-\boldsymbol\tau^{\otimes N}\|_1
                        \le d_N(\mathcal R_\sigma,\mathcal C).
\]
Lemma~\ref{lem:productreplacer68} identifies the left side with the left
side of \eqref{eq:centredominance70}. Retracting every centre in any cover
preserves coverage and cannot increase cardinality. The reverse inequality
follows because preparation centres are among the legal centres. The same
argument preserves rationality.
\end{proof}

The retraction does not preserve a target rank bound: for example the
identity channel retracts to the maximally mixed preparation. Thus the
centre restriction in Theorem~\ref{thm:retraction70} is to the full
preparation family, not necessarily to the target rank stratum. The
rank-constrained rational upper centres of
Theorem~\ref{thm:factorcode68} remain a separate, stronger construction.
The theorem concerns the geometry of mathematical descriptions. Its
maximally mixed test is a quantum experiment, not a finite-classical-message
simulation of unknown inputs.

The reduction mechanism for input-erasing channels has a direct antecedent
in Cooney--Mosonyi--Wilde \cite{CMW70}. Their discussion of two replacer
channels reduces channel divergences to state divergences; their principal
results concern asymptotic discrimination and strong-converse exponents
for a channel versus a replacer. Lemma~\ref{lem:productreplacer68} uses the
same input-erasure mechanism in an exact finite-use trace-norm identity,
whereas Theorem~\ref{thm:retraction70} concerns the allowed centres of an
entire covering problem. No exact nonadaptive optimality for an arbitrary
non-replacer channel against a replacer is inferred from that identity.
The processor in the product identity is quantum, and input erasure,
rather than entanglement breaking alone, is its essential hypothesis.
